% Standalone implementation plan.
\documentclass[11pt,a4paper]{article}
\usepackage[T1]{fontenc}
\usepackage[utf8]{inputenc}
\usepackage{lmodern}
\usepackage[margin=25mm]{geometry}
\usepackage{amsmath,amssymb}
\usepackage{booktabs}
\usepackage{array}
\usepackage{microtype}
\usepackage[numbers,sort&compress]{natbib}
\usepackage{xurl}
\usepackage[hidelinks]{hyperref}

\newcommand{\R}{\mathbb R}
\newcommand{\col}{\operatorname{col}}
\newcolumntype{L}[1]{>{\raggedright\arraybackslash}p{#1}}

\title{Implementation plan for one-step orthogonal augmentation of the gantry model}
\author{Eindhoven University of Technology}
\date{16 September 2026}

\begin{document}
\maketitle

\section{Objective and scope}

The first implementation imposes orthogonality in one evaluation of the
gantry plant state equation. It retains the physical-plus-augmentation state
structure of Hoekstra et al.~\cite{hoekstra2025lfr,hoekstra2026fp}; no
delay-coordinate state is introduced. The projection and its training
semantics follow Gy\"or\"ok et al.~\cite{gyorok2026obc}. The local
parameter-sensitivity construction is informed by the earlier state-space
regularization method of Gy\"or\"ok et al.~\cite{gyorok2025l4dc}.

The protected space is the tangent space generated by the ten identifiable
baseline parameter combinations. The construction does not impose
orthogonality of a multistep input--output or closed-loop Jacobian. The
controller and recursive propagation remain active in training, but their
derivatives do not enter the first orthogonality condition.

\section{Source methods and implementation starting point}

Inspect the two local source repositories before implementation.
\begin{enumerate}
    \item \path{orthogonal-IO-augm-main} is the authority for the
    orthogonal-by-construction subtraction. In
    \path{orthogonalx_augm/src.py}, class
    \path{IO_orthogonal_augmentation}, the training regressor matrix and its
    factorization are constructed outside the objective, while the auxiliary
    coefficient is recomputed inside every objective evaluation:
    \begin{equation}
        \theta_{\mathrm{aux}}=K_\Phi F_{\mathrm{ANN}},
        \quad K_\Phi=(\Phi^\top\Phi)^\dagger\Phi^\top,
        \quad \widetilde F_{\mathrm{ANN}}
        =F_{\mathrm{ANN}}-\Phi\theta_{\mathrm{aux}}.
        \label{eq:source-code-obc}
    \end{equation}
    Gradients pass through \(\theta_{\mathrm{aux}}\) to the ANN. After
    training, the final coefficient is frozen and the regressor is evaluated
    at each prediction point, as in Gy\"or\"ok's Eq.~(13).

    \item \path{orthogonal-augmentation-main} is the state-space predecessor.
    It supplies sample-major stacking and SVD conventions and motivates the
    parameter Taylor expansion in Eqs.~(15)--(19) of
    \cite{gyorok2025l4dc}. It implements soft regularization and does not
    implement dynamic additional states, so its penalty does not define the
    present by-construction method.
\end{enumerate}

Use \path{Reduced_Gantry_State_Block} in
\path{model_augmentation/fit_systems/blocks.py}. The OBC experiment requires
\begin{center}
    \texttt{joint\_estimation=True},\qquad
    \texttt{physics\_parameterization='reduced'}.
\end{center}
The general training entry point currently has
\texttt{joint\_estimation=False}; the OBC comparison must enable it explicitly.
The block trains
\begin{equation}
\vartheta_b=\col(
 k_{b,\mathrm{sum}},c_{g1},c_{g2},c_y,c_{b,\mathrm{sum}},m_h,
 m_{\mathrm{total}},m_{\mathrm{diff}},J_{\mathrm{eff}},d)
 \in\R^{10}.
\label{eq:implemented-reduced-parameters}
\end{equation}
The raw fourteen parameters reconstructed through the fixed gauge are not
estimated coordinates. Apply the initial perturbation directly to these ten
combinations through \path{combo_init_detune}. The intended factors, in the
order of Eq.~\eqref{eq:implemented-reduced-parameters}, are
\begin{equation}
    (1.10,1.10,0.90,1.10,0.90,0.90,1.10,1.10,0.90,1.10).
\end{equation}
Differentiate the production block's complete discrete-time RK4 transition;
do not duplicate the gantry dynamics in the OBC implementation.

\section{One-step augmented model}

Let
\begin{equation}
    \hat x_k=\col(x_{b,k},x_{a,k}),
    \qquad x_{b,k}\in\R^{n_b},\quad x_{a,k}\in\R^{n_a}.
\end{equation}
The augmented model is
\begin{equation}
\boxed{
\begin{aligned}
    x_{b,k+1}
    &=f_{\mathrm{base}}^d(\vartheta_b,x_{b,k},u_k)
      +f_{\mathrm{aug}}^b(\theta_a,\hat x_k,u_k),\\
    x_{a,k+1}
    &=f_{\mathrm{aug}}^a(\theta_a,\hat x_k,u_k),\\
    \hat y_k&=h_{\mathrm{base}}(x_{b,k},u_k).
\end{aligned}}
\label{eq:model}
\end{equation}
For the current gantry configuration, \(n_b=6\), \(n_a=8\), and the complete
augmented state has dimension 14.

The baseline has no transition contribution in the additional-state rows.
For a baseline parameter direction \(\phi_{b,k}\),
\begin{equation}
    \bar\phi_{b,k}=\begin{bmatrix}\phi_{b,k}\\0\end{bmatrix},
    \qquad
    \bar\phi_{b,k}^{\top}
    \begin{bmatrix}0\\f_{\mathrm{aug},k}^a\end{bmatrix}=0.
    \label{eq:zero-block}
\end{equation}
Thus, at the one-step augmented-state level, the additional-state rows are
orthogonal to the baseline by the zero-block structure and are left
unconstrained. This does not constrain how nonzero \(x_a\) later influences
\(f_{\mathrm{aug}}^b\), nor its accumulated influence through recursion.

\section{Fixed one-step reference set}
\label{sec:fixed-reference-set}

Use every valid sample in the complete training dataset. The measured output
contains the three stage positions and satisfies
\begin{equation}
    y_k=P^\top q_k.
\end{equation}
Reconstruct the logical positions and estimate the unmeasured velocities by
\begin{equation}
    q_k^{\mathrm{ref}}=P^{-\top}y_k^{\mathrm{data}},
    \qquad
    \widehat{\dot q}_k^{\mathrm{FD}}
    =\frac{-q_{k+2}^{\mathrm{ref}}+8q_{k+1}^{\mathrm{ref}}
            -8q_{k-1}^{\mathrm{ref}}+q_{k-2}^{\mathrm{ref}}}{12T_s},
    \label{eq:reference-reconstruction}
\end{equation}
using suitable lower-order formulas at record boundaries. Position
reconstruction is an exact coordinate inversion of the implemented output
map. The finite difference is a data-derived velocity estimate, not an exact
state measurement. It is affected by noise and frequency-dependent gain. The
fourth-order stencil limits the latter in the active 140--230 Hz band: at
4 kHz its gain error at 230 Hz is below 0.2\%, compared with approximately
2.2\% for the second-order central difference. This matters because the
damping directions multiply velocity.

Define
\begin{equation}
    x_{b,k}^{\mathrm{ref}}
    =\col(q_k^{\mathrm{ref}},\widehat{\dot q}_k^{\mathrm{FD}}),
    \qquad x_{a,k}^{\mathrm{ref}}=0,
    \qquad u_k^{\mathrm{ref}}=u_k^{\mathrm{data}}.
    \label{eq:reference-state}
\end{equation}
The recorded input is the plant input at the measured point. It is held fixed
when differentiating the plant; no controller derivative is included. The
fixed reference set is
\begin{equation}
\boxed{
    \mathcal Z_{\mathrm{ref}}
    =\left\{(x_{b,k}^{\mathrm{ref}},0,u_k^{\mathrm{data}})
      :k\in\mathcal I_{\mathrm{train}}\right\}.
}
\label{eq:reference-set}
\end{equation}
Its elements are independent one-step evaluations, not a recursively simulated
trajectory. The actual training rollout remains closed loop and may continue
to train the encoder.

The choice \(x_a^{\mathrm{ref}}=0\) defines the first condition on the slice
\(\{x_a=0\}\). Thus, \(f_{\mathrm{aug}}^b(x_b,x_a,u)\) is constrained on
that slice only. Directions involving nonzero additional states remain outside
the first condition.

\section{Baseline tangent space}

Gy\"or\"ok's input--output construction has the exact representation
\begin{equation}
    f_{\mathrm{base}}(z_k;\theta_b)=\phi(z_k)\theta_b.
\end{equation}
Its fixed stacked regressor \(\Phi\) spans the baseline output directions. The
nonlinear gantry transition has no exact parameter-independent regressor
because its parameters enter the inverse mass matrix and RK4 stages
nonlinearly.

At reference point \(k\), use
\begin{equation}
    J_{b,k}
    =\left.
      \frac{\partial f_{\mathrm{base}}^d
      (\vartheta_b,x_{b,k}^{\mathrm{ref}},u_k^{\mathrm{data}})}
      {\partial\vartheta_b}
      \right|_{\vartheta_b=\bar\vartheta_b}
    \in\R^{n_b\times10}.
    \label{eq:local-sensitivity}
\end{equation}
This gives the local expansion
\begin{equation}
 f_{\mathrm{base}}^d(\bar\vartheta_b+\Delta\vartheta_b,x_b,u)
 \approx f_{\mathrm{base}}^d(\bar\vartheta_b,x_b,u)
 +J_b(\bar\vartheta_b,x_b,u)\Delta\vartheta_b.
\end{equation}
The primary protected space is
\begin{equation}
    \mathcal S_b=\operatorname{range}(J_b),
    \label{eq:protected-space}
\end{equation}
the stacked first-order transition changes produced by the ten identifiable
parameter combinations. This choice targets parameter--ANN overlap.

The nonlinear state-space construction in Eqs.~(17)--(19) of
\cite{gyorok2025l4dc} instead protects
\begin{equation}
    \operatorname{range}([J_b\mid\Gamma]),
    \qquad
    \Gamma_k=f_{\mathrm{base}}^d(\bar\vartheta_b,z_k)
              -J_{b,k}\bar\vartheta_b.
    \label{eq:extended-space}
\end{equation}
This also contains the affine offset. It is not the primary choice because a
discrete transition is dominated by identity propagation, so protecting
\(\Gamma\) can exclude a broad state-like ANN direction unrelated to parameter
ambiguity. The extended space is the first follow-up comparison.

Existing investigation results used the extended eleven-column space
\([J_b\mid\Gamma]\), not the ten-column primary space. They therefore neither
establish benefit nor establish harm for Eq.~\eqref{eq:protected-space}.
Before training, evaluate the discrepancy-overlap diagnostic for both spaces:
\begin{equation}
    \rho(\mathcal B;\delta)
    =\frac{\|\mathcal B\mathcal B^\dagger\delta\|_2}
           {\|\delta\|_2+\varepsilon},
    \qquad
    \mathcal B\in\{J_b,[J_b\mid\Gamma]\}.
    \label{eq:discrepancy-overlap}
\end{equation}
Use the same reference points and discrepancy in both calculations. This is a
required diagnostic only when a defensible one-step discrepancy \(\delta\) is
available. A large value shows that Gy\"or\"ok's recovery condition is
unfavourable for that protected space; it does not by itself invalidate the
projection or determine its prediction performance.

\section{Scheduling variable}

Use the scheduling value at each reference point:
\begin{equation}
    Y_k^{\mathrm{ref}}
    =[x_{b,k}^{\mathrm{ref}}]_3=[y_k^{\mathrm{data}}]_3.
\end{equation}
Evaluate the production LPV transition and sensitivity at this exact value. No
Taylor expansion in \(Y\) is used. Inspect coverage of the training range of
\(Y\), because insufficient variation can make scheduling-dependent parameter
directions poorly conditioned.

\section{Coordinates, routing and stacking}

The production physical block already accepts normalized state and input
coordinates and returns the normalized next state. The ANN writes into the
same normalized rows. Autodifferentiating the block therefore gives
\begin{equation}
    J_{b,k}^{\mathrm n}
    =\left.
      \frac{\partial f_{\mathrm{base}}^{d,\mathrm n}}
      {\partial\vartheta_b}
      \right|_{z_k^{\mathrm{ref}},\bar\vartheta_b}.
    \label{eq:normalized-sensitivity}
\end{equation}
Do not apply \(S_b^{-1}\) again to this result or the ANN output. Euclidean
orthogonality in normalized state coordinates corresponds to a
variance-weighted metric in physical coordinates. It is distinct from the
closed-loop rollout loss on normalized outputs.

Let \(R_b\in\R^{n_r\times n_b}\) select the physical rows that the ANN can
write. Use the same selection and order for the baseline and ANN:
\begin{equation}
    \mathcal J_k=R_bJ_{b,k}^{\mathrm n}\in\R^{n_r\times10},
    \qquad
    g_k=R_bf_{\mathrm{aug}}^{b,\mathrm n}(z_k^{\mathrm{ref}})
       \in\R^{n_r}.
\end{equation}
Stack sample-major over all \(N\) reference samples:
\begin{equation}
    J_b^{\mathrm n}
    =\begin{bmatrix}\mathcal J_1\\ \vdots\\ \mathcal J_N\end{bmatrix}
       \in\R^{Nn_r\times10},
    \qquad
    F_{\mathrm{aug}}^{b,\mathrm n}
    =\begin{bmatrix}g_1\\ \vdots\\ g_N\end{bmatrix}
       \in\R^{Nn_r}.
    \label{eq:stacks}
\end{equation}
This is a dataset stack, not a product of transition Jacobians. Require
\begin{equation}
    \operatorname{rank}(J_b^{\mathrm n})=10.
    \label{eq:full-column-rank}
\end{equation}
Full column rank gives a unique least-squares coefficient. The projected vector
is still defined by the Moore--Penrose pseudoinverse without full rank, but a
rank loss means that the experiment does not separate all ten selected
parameter directions and invalidates the intended recovery test. If the rank
is below ten, determine whether the selected combinations are still redundant
or whether the training experiment lacks excitation, and correct that issue
before training. Do not silently continue with a rank-truncated version of the
intended ten-parameter recovery experiment.

For the current 14-record dataset, each record contains approximately
48,000 usable samples. Hence \(N\) is approximately \(6.7\times10^5\). With
all six physical rows routed, \(n_r=6\), so the complete stack has roughly
\(4.0\times10^6\) rows and ten columns. Recompute and log the exact values from
the loaded data rather than relying on these planning estimates.

\section{Projection and training lifecycle}

Construct the pseudoinverse from a direct economy SVD,
\begin{equation}
    J_b^{\mathrm n}=U\Sigma V^\top,
    \qquad K_J=(J_b^{\mathrm n})^\dagger=V\Sigma^\dagger U^\top,
    \label{eq:svd}
\end{equation}
with a declared rank tolerance. Do not form and invert the Gram matrix, which
squares the condition number.

Construct and factorize the reference basis in \texttt{float64}, including the
rank and residual diagnostics. Cast the current ANN reference field to
\texttt{float64} when calculating \(\theta_{\mathrm{aux}}\), preserving its
autograd connection, and cast the resulting coefficient to the pipeline dtype
only at the rollout boundary. Report both dtypes. Judge the numerical-zero
test against a declared scale such as
\begin{equation}
    \tau_\perp=c\,\kappa(J_b^{\mathrm n})
    \epsilon_{\mathrm{mach}},
\end{equation}
with the constant \(c\), condition number and machine precision recorded,
rather than using an unspecified ``near zero'' threshold.

For the current ANN output on the reference set, define
\begin{equation}
\begin{aligned}
    \theta_{\mathrm{aux}}
    &=K_JF_{\mathrm{aug}}^{b,\mathrm n}\in\R^{10},\\
    \widetilde F_{\mathrm{aug}}^{b,\mathrm n}
    &=F_{\mathrm{aug}}^{b,\mathrm n}
      -J_b^{\mathrm n}\theta_{\mathrm{aux}}.
\end{aligned}
\label{eq:projection}
\end{equation}
Then
\begin{equation}
    \boxed{(J_b^{\mathrm n})^\top
    \widetilde F_{\mathrm{aug}}^{b,\mathrm n}=0}
    \label{eq:orthogonality}
\end{equation}
to numerical precision. This is one stacked condition over the reference set,
not a pointwise condition.

\subsection{Epoch boundary}

At the start of epoch \(r\), set
\begin{equation}
    \bar\vartheta_b^{(r)}=\widehat\vartheta_b^{(r)}.
\end{equation}
Build \(J_b^{\mathrm n,(r)}\), report its complete singular-value spectrum,
rank and condition number, and construct
\begin{equation}
    K_J^{(r)}=(J_b^{\mathrm n,(r)})^\dagger.
\end{equation}
Store the basis, pseudoinverse and expansion point as detached buffers for the
epoch. No gradient passes through basis construction. If the physical
parameters are fixed, build the basis once because its expansion point does
not move.

\subsection{Every objective evaluation}

Evaluate the current ANN once over the complete \(\mathcal Z_{\mathrm{ref}}\)
and compute
\begin{equation}
    \theta_{\mathrm{aux}}(\theta_a)
    =K_J^{(r)}F_{\mathrm{aug}}^{b,\mathrm n}(\theta_a).
    \label{eq:coefficient-per-objective}
\end{equation}
Do not detach or cache this coefficient across optimizer updates. Gradients
must pass through
\(\theta_a\mapsto F_{\mathrm{aug}}^{b,\mathrm n}\mapsto
\theta_{\mathrm{aux}}\). This matches Gy\"or\"ok: his \(\Phi\) and
\(K_\Phi\) are fixed, while \(\theta_{\mathrm{aux}}\) is recomputed inside
every objective evaluation.

Compute \(\theta_{\mathrm{aux}}\) once for the objective and reuse it at every
timestep of that rollout. Re-evaluating the complete reference set inside
every model timestep is redundant and prohibitively expensive. At rollout
point \((\hat x_k,u_k)\), apply
\begin{equation}
    \widetilde f_{\mathrm{aug}}^{b,\mathrm n}(\hat x_k,u_k)
    =f_{\mathrm{aug}}^{b,\mathrm n}(\hat x_k,u_k)
     -J_b^{\mathrm n}(\bar\vartheta_b^{(r)},x_{b,k},u_k)
      \theta_{\mathrm{aux}}.
    \label{eq:pointwise-correction}
\end{equation}
Evaluate \(J_b\theta_{\mathrm{aux}}\) as one forward-mode JVP of the
production transition in direction \(\theta_{\mathrm{aux}}\). Leave
\(f_{\mathrm{aug}}^a\) unchanged.

The SVD remains outside the transition, but the JVP adds work at every rollout
timestep. Benchmark its runtime and memory and its interaction with gradient
checkpointing, \path{torch.compile}, and the physical block's \path{_cur}
cache before a full run.

The complete training dataset defines \(\mathcal Z_{\mathrm{ref}}\) in the
primary method. If evaluating it during every objective proves infeasible,
introduce a decimated auxiliary set only as a measured approximation. Compare
its rank, protected subspace and overlap metrics with the complete-data result.

\section{Required experiments and diagnostics}

Use identical data, initialization, controller, normalization, optimizer and
random seed for:
\begin{enumerate}
    \item the jointly estimated augmented model without OBC;
    \item the jointly estimated augmented model with one-step OBC.
\end{enumerate}
Initialize the same ten reduced combinations in both runs. Also report the
following two FP references separately:
\begin{itemize}
    \item the open-loop recorded-input FP simulation as a plant-replay
          diagnostic;
    \item the controller-wrapped FP simulation with the ANN disabled, using
          the same closed-loop scoring convention as the trained models.
\end{itemize}
If the supervisor's ``regular orthogonalization method'' refers to the
published 2025 soft regularization rather than the one-step level of the 2026
construction, add a third trained arm using the existing 2025 penalty. This
arm is also required when a direct penalty-versus-subtraction contribution
claim is made. Otherwise, keep it conditional to limit the run count.

At every basis refresh, report the singular values, numerical rank and
condition number. During training, report
\begin{equation}
    r_\perp=
    \frac{\|(J_b^{\mathrm n})^\top
    \widetilde F_{\mathrm{aug}}^{b,\mathrm n}\|_2}
    {\|J_b^{\mathrm n}\|_F
     \|\widetilde F_{\mathrm{aug}}^{b,\mathrm n}\|_2+\varepsilon}
\end{equation}
and
\begin{equation}
    r_{\mathrm{overlap}}
    =\frac{\|J_b^{\mathrm n}(J_b^{\mathrm n})^\dagger
    F_{\mathrm{aug}}^{b,\mathrm n}\|_2}
    {\|F_{\mathrm{aug}}^{b,\mathrm n}\|_2+\varepsilon}.
\end{equation}
With Eq.~\eqref{eq:coefficient-per-objective}, \(r_\perp\) must remain below
the declared dtype- and conditioning-dependent tolerance at every checked
objective evaluation.

Monitor movement of the baseline tangent:
\begin{equation}
    r_J^{(r)}
    =\frac{\|J_b^{\mathrm n,(r)}-J_b^{\mathrm n,(0)}\|_F}
    {\|J_b^{\mathrm n,(0)}\|_F+\varepsilon},
    \qquad
    r_{J,\mathrm{step}}^{(r)}
    =\frac{\|J_b^{\mathrm n,(r)}-J_b^{\mathrm n,(r-1)}\|_F}
    {\|J_b^{\mathrm n,(r-1)}\|_F+\varepsilon}.
\end{equation}
Report each reduced-parameter error, the aggregate error
\begin{equation}
    e_{\vartheta}
    =\frac{\|\widehat\vartheta_b-\vartheta_b^{\mathrm{true}}\|_2}
    {\|\vartheta_b^{\mathrm{true}}\|_2},
\end{equation}
one-step prediction error, and closed-loop training, validation and free-run
errors. When the data-generating discrepancy is available, report its overlap
with the protected space because Gy\"or\"ok's recovery result assumes that the
unmodelled term is orthogonal to the baseline regressor.

\section{Decision criteria}

Interpret the first results using the following pre-committed criteria.
\begin{center}
\small
\begin{tabular}{L{0.29\linewidth}L{0.34\linewidth}L{0.27\linewidth}}
\toprule
Observation & Interpretation & Action \\
\midrule
Rank below ten & The chosen combinations remain redundant or the data do not
excite every direction & Resolve the combinations or experiment; do not start
the recovery comparison \\
\(r_\perp>\tau_\perp\) & The by-construction implementation is incorrect or
numerically inadequate & Debug before interpreting training results \\
Large \(\rho(J_b;\delta)\) & The discrepancy violates the recovery condition
for the primary tangent space & Do not claim expected parameter recovery;
compare with the no-OBC arm \\
Small \(\rho(J_b;\delta)\), large
\(\rho([J_b\mid\Gamma];\delta)\) & The affine offset causes the overlap & Keep
\(J_b\) primary and treat the extended space as a diagnostic \\
Large \(r_J^{(r)}\) & The initial tangent becomes stale & Keep the epoch-wise
basis refresh; use the fixed-basis arm only as a diagnostic \\
Orthogonality holds but recovery does not improve & The one-step space,
reference states, or recovery assumptions are insufficient & Inspect
discrepancy overlap and reference-state sensitivity before multistep methods \\
Recovery improves but closed-loop error does not & Overlap reappears through
propagation or is unrelated to predictive error & Investigate multistep or
closed-loop sensitivities \\
Complete-data reference evaluation is infeasible & The exact chosen reference
definition is too expensive & Validate a reduced auxiliary set before use \\
\bottomrule
\end{tabular}
\end{center}

\section{Conditional follow-up experiments}

Run these only when the primary result identifies the corresponding need.
\begin{enumerate}
    \item \textbf{Extended affine space.} Replace \(J_b\) by
    \([J_b\mid\Gamma]\) from Eq.~\eqref{eq:extended-space}.
    \item \textbf{Fixed initial basis.} Retain
    \(J_b^{\mathrm n,(0)}\), but continue to recompute
    \(\theta_{\mathrm{aux}}\) during every objective evaluation.
    \item \textbf{Oracle parameter point.} Use
    \(\bar\vartheta_b=\vartheta_b^{\mathrm{true}}\) for diagnosis only.
    \item \textbf{Alternative velocity reference.} Compare finite-difference
    velocities with a fixed encoder or baseline-model estimate if noise makes
    the derivative unreliable.
    \item \textbf{Reference-set reduction.} Introduce a declared stride only
    if the complete-data objective is infeasible; compare singular values,
    principal angles and overlap metrics first.
    \item \textbf{Nonzero additional-state points.} Sample plausible \(x_a\)
    values if the zero slice proves insufficient.
    \item \textbf{Multistep or closed-loop orthogonality.} Introduce products
    of plant, output and controller Jacobians only if the one-step construction
    works while overlap reappears through propagation.
\end{enumerate}

\section{Claims after implementation}

The implementation may claim that:
\begin{itemize}
    \item the additional-state rows are orthogonal to the baseline at the
          one-step augmented-state level because the baseline block is zero in
          those rows;
    \item the corrected physical ANN contribution is orthogonal, over the
          complete reference stack, to the selected baseline parameter tangent
          at the current epoch's expansion point;
    \item this condition is exact to numerical precision for the current ANN
          at each objective evaluation and local with respect to the nonlinear
          baseline parameter family;
    \item no pointwise, nonzero-\(x_a\), multistep, input--output, stability or
          closed-loop orthogonality follows from this construction.
\end{itemize}
Gy\"or\"ok's parameter-recovery, consistency, covariance and stability results
do not transfer automatically to the nonlinear, marginally stable, closed-loop
gantry model.

\bibliographystyle{plainnat}
\bibliography{obc-io-additional-state-chapter}

\end{document}
